% Fuentes: tab_s09_errores_predictivos.csv
\begin{table}[htbp]
\centering
\caption{Rango de RMSE en tres bloques internos y error del bloque final a $1200$ s.}
\label{tab:res_validacion_ml}
\small
\setlength{\tabcolsep}{5pt}
\begin{tabular}{lrrrr}
\toprule
Variable & Modelo & Mínimo interno & Máximo interno & Prueba final \\
\midrule
$T_{\mathrm{DS}}$ & XGBoost & $0{,}1765$ & $0{,}3973$ & $0{,}1765$ \\
$T_{\mathrm{DS}}$ & LightGBM & $0{,}1765$ & $0{,}3774$ & $0{,}1758$ \\
$T_{\mathrm{DHT}}$ & XGBoost & $0{,}2019$ & $0{,}4291$ & $0{,}2014$ \\
$T_{\mathrm{DHT}}$ & LightGBM & $0{,}2016$ & $0{,}4097$ & $0{,}2020$ \\
$\phi_{\mathrm{DHT}}$ & XGBoost & $1{,}0496$ & $1{,}4591$ & $1{,}0622$ \\
$\phi_{\mathrm{DHT}}$ & LightGBM & $1{,}0445$ & $1{,}4995$ & $1{,}0647$ \\
\bottomrule
\end{tabular}
\par\vspace{3pt}
\begin{minipage}{0.98\textwidth}\footnotesize
Los extremos corresponden a bloques de origen rodante. El rango no es un intervalo de confianza y no se transforma en una media agrupada sin los errores individuales y sus pesos.
\end{minipage}
\end{table}
